\documentclass[book.tex]{subfiles}

\begin{document}

\chapter{Recurrent Neural Networks}

\label{chap:rnns}
\noindent\rule{11cm}{0.4pt} \\
\textbf{Prerequisites:}  \autoref{chap:ml_basics}, \autoref{chap:grad_descent}  \\
\textbf{Difficulty Level:} ** \\
\noindent\rule{11cm}{0.4pt}
\vspace{5mm}

Recurrent neural networks (RNNs) make use of sequential structure in a data source. RNNs learn a transformation that is repeated at each  step to process input data. Usually, the sequential structure corresponds to time, but RNNs have also been applied to sequential data such as natural language texts or genomic sequences. 
\begin{figure}
    % \centering    \includegraphics{figures/Differentiable Models/rnns/Recurrent_neural_network_unfold.svg.png}
   \centering \includegraphics{figures/Differentiable Models/rnns/Self Connection-neural network1.png}
    \caption{A recurrent neural network unfolds a self-connection out over time. }
    \label{fig:my_label}
\end{figure}

\section{Recurrent Neural Network Equations}

Assume that we have a series of inputs $x_1,\dotsc,x_T$. A recurrent network constructs a series of hidden state vectors $h_1,\dotsc, h_T$ for each time step in the network. These hidden states are updated based at each time step $t$ based on the input $x_t$ and the previous hidden state $h_{t-1}$. The update rule is parameterized by a set of weights $\theta$. Formally, the update equations are

\begin{align}
\hat{y}_t, h_t = f(h_{t-1}, x_t; \theta)
\end{align}
Here $f$ is some transformation function, commonly implemented by a fully connected network. 
The total loss is the sum of the losses at each timestep for each output prediction.
\begin{align}
    \mathcal{L}_{\textrm{total}} = \sum_{i=1}^T \mathcal{L}(\hat{y}_t, y_t)
\end{align}

The gradient update for all parameters $\Theta$ is computed as usual with a gradient descent step.
\begin{align}
    \Theta = \Theta - \alpha \frac{\partial \mathcal{L}_{\textrm{total}}}{\partial \Theta}
\end{align}

%We define softmax as follows
%\begin{lstlisting}
%def softmax(x: $\mathbb{R}^n$) $\to$ $\mathbb{R}^n$:
%  p = exp(x - max(x))
%  return p / sum(p)
%\end{lstlisting}
Here is the code for the full model definition. %\textcolor{red}{TODO: 

\begin{lstlisting}[language=python]
class RNN(f: FullyConnectedNetwork, N : $\mathbb{N}$, d : $\mathbb{N}$, T : $\mathbb{N}$):
  def $\lambda$(X : $\mathbb{R}$[N, T, d]) $\to$ $\mathbb{R}$[N, T]:
    out = zeros(N, T)
    for i in N:
      x, y = X[i], Y[i] 
      ht = zeros(H, 1)
      for t in T:
        out[i], ht = f(ht || x[t])
    out
\end{lstlisting}

Here we use the operator \lstinline{||} to denote concatenation of vectors. 



\section{Vanishing Gradients and The Challenge of Long$-$Term Dependencies}

Recurrent neural networks often struggle to handle very long input sequences. Since the loss function $\mathcal{L}_{\textrm{total}}$ includes terms from all time steps, gradients will have to propagate backwards through long chains of updates. Each gradient update step can be viewed conceptually as adding a little bit of noise to a propagating signal, which means that gradients become less useful deeper back in time. For this reason, many RNN implementations truncate backwards updates in time past some set number of steps (say 50 or 100).

The basic problem is that gradients propagated over many stages tend to either vanish (most of the time) or explode (rarely, but with much damage to the optimization). Even it is assumed that the parameters are such that the recurrent network is stable (can store memories, with gradients not exploding), the difficulty with long-term dependencies arises from the exponentially smaller weights given to long-term interactions (involving the multiplication of many Jacobians) compared to short-term ones. Recurrent networks involve the composition of the same function multiple times, once per time step. These compositions can result in extremely nonlinear behavior. In particular, the function composition employed by recurrent neural networks somewhat resembles matrix multiplication. Take the equation
\begin{align}
h^{(t)} &= W^Th^{(t-1)}
\end{align}
as a very simple recurrent neural network lacking a nonlinear activation function, and lacking inputs x. It may be simplified to 
\begin{align}
    h^{(t)} = (W^t)^Th^{(0)}
\end{align}
and if W admits an eigendecompostion of the form
\begin{align}
    W=Q\Lambda Q^T
\end{align}
with orthogonal Q, the recurrence may be simplified further to
\begin{align}
    h^{(t)}=Q^T\Lambda^tQh^{(0)}
\end{align}
The eigenvalues are raised to the power of t causing eigenvalues with magnitude less than one to decay to zero and eigenvalues with magnitude greater than one to explode. Any component of $h^{(0)}$ that is not aligned with the largest eigenvector will eventually be discarded. This problem remains an active area of research within machine learning. 

\section{More Sophisticated Update Rules}

In the discussion above, we introduced the time update function $f$ as a fully connected network, but more complex update rules have been found empirically to improve performance, especially with respect to long timescale behavior. Here is one such update rule that uses the $\mathrm{tanh}$ function for example. 
\begin{align}
    a_t &= b + W h_{t-1} + U x_{t} \\
    h_t &= \tanh(a_t) \\
    o_t &= c + V h_{t} \\
    \hat{y}_t &= \textrm{softmax}(o_t)
\end{align}

Here are a few other structural changes which can help improve RNN performance

\begin{enumerate}
    %\item Echo State Networks: set the recurrent weights such that the recurrent hidden units do a good job of capturing the history of past inputs, and learn only the output weights.
    \item Use time skips so that rather than going from $t$ to $t+1$, the network goes from $t$ to $t+d$. This allows the gradients to diminish much more slowly.
    \item Use units with linear self-connections and a weight near one on these connections. Accumulate a running average, $\mu^{(t)}$, of some value $v^{(t)}$ by applying the update equation 
    \begin{align}\mu^{(t)} = \mu^{(t-1)}\alpha + (1-\alpha)v^{(t)}\end{align}
    where $\alpha$ is the linear self-connection and if it is near one, the running average remembers information about the past for a long time, and when it is near zero, information about the past is rapidly discarded.
    %\item Adjust connections of the neural network such that they operate over longer time scales
    \item Leverage LSTM (Long Short$-$Term Memory) structure. Currently the most effective sequence model and used in most RNNs. This model is a gated RNN which allows the network to accumulate information over a long duration and then decide whether or not to forget some of that information once it has been used. The LSTM includes intermediate steps such as "input gates", "output gates" and "forget gates" that can help preserve internal structure. %The LSTM model is shown in the margin.
\end{enumerate}

\begin{figure}
     % \centering    \includegraphics{figures/Differentiable Models/rnns/lstm_cell.png}
    \centering    \includegraphics{figures/Differentiable Models/rnns/Internal Structure of LSTM Cell.png}
    \caption{The internal structure of an LSTM cell.}
    \label{fig:my_label}
\end{figure}

\end{document}